\chapter{Number theory}

\section{Modular arithmetic}
	\kactlimport{ModularArithmetic.h}
	\kactlimport{ModInverse.h}
	\kactlimport{ModPow.h}
	\kactlimport{ModLog.h}
	\kactlimport{ModSum.h}
	\kactlimport{ModMulLL.h}
	\kactlimport{ModSqrt.h}
	\kactlimport{PrimitiveRoot.h}

\section{Primality}
	\kactlimport{FastEratosthenes.h}
	% \kactlimport{LinearSieve.h}
	\kactlimport{MillerRabin.h}
	\kactlimport{Factor.h}

\section{Divisibility}
	\kactlimport{euclid.h}
	% \kactlimport{Euclid.java}
	\kactlimport{CRT.h}
	\kactlimport{ExLucas.h}

	\subsection{Bézout's identity}
	For $a \neq 0$, $b \neq 0$, then $d=gcd(a,b)$ is the smallest positive integer for which there are integer solutions to
	$$ax+by=d$$
	If $(x,y)$ is one solution, then all solutions are given by
	$$\left(x+\frac{kb}{\gcd(a,b)}, y-\frac{ka}{\gcd(a,b)}\right), \quad k\in\mathbb{Z}$$
	$ax+by=c$ is solvable iff $g=\gcd(a,b) \mid c$: take $(x_0,y_0) = (x_g, y_g) \cdot c/g$ from \texttt{euclid}, then $x = x_0 + \frac bg t$, $y = y_0 - \frac ag t$.
	For $a, b > 0$, $x \in [L, R]$ means $t \in [\lceil (L-x_0)g/b \rceil, \lfloor (R-x_0)g/b \rfloor]$ and $y \in [L', R']$ means $t \in [\lceil (y_0-R')g/a \rceil, \lfloor (y_0-L')g/a \rfloor]$:
	intersect the ranges (careful with floor division of negatives). $x + y$ changes by $(b-a)/g$ per step of $t$, so it is minimal at an end of the range.

	\kactlimport{phiFunction.h}

	\subsection{Garner's algorithm}
		% Checked by doc/note-checks/check_numbertheory.py
		CRT for pairwise coprime $m_0, \dots, m_{k-1}$ (for two moduli see \texttt{crt}): write $x = x_0 + x_1 m_0 + x_2 m_0 m_1 + \dots + x_{k-1} m_0 \cdots m_{k-2}$ with $0 \le x_i < m_i$. Then
		$x_i = \bigl(a_i - x_0 - x_1 m_0 - \dots - x_{i-1} m_0 \cdots m_{i-2}\bigr)(m_0 \cdots m_{i-1})^{-1} \bmod m_i$, in $O(k^2)$ with the inverses $m_j^{-1} \bmod m_i$ precomputed:\\
		\verb@for i: t = a[i];@\\
		\verb@  for j < i: t = (t-x[j])@\\
		\verb@    * inv(m[j], m[i]) % m[i];@\\
		\verb@  x[i] = t;@\\
		Then $x \bmod M$ for any $M$ by Horner (e.g.\ a convolution result mod $M$ from three NTT primes):\\
		\verb@r = x[k-1]; for i = k-2..0:@\\
		\verb@  r = (r*m[i] + x[i]) % M;@

	\subsection{Exponents}
		Generalized Euler: $a^b \equiv a^{(b \bmod \phi(m)) + \phi(m)} \pmod m$ for every $a$ and every $b \ge \log_2 m$ (use it for $a^{b^c}$ and power towers).
		Lifting the exponent: for an odd prime $p$ with $p \mid x - y$ and $p \nmid xy$, $v_p(x^n - y^n) = v_p(x-y) + v_p(n)$;
		if $n$ is odd and $p \mid x + y$, $v_p(x^n + y^n) = v_p(x+y) + v_p(n)$.
		For $p = 2$ and odd $x, y$: $v_2(x^n - y^n) = v_2(x-y)$ for odd $n$, and $v_2(x-y) + v_2(x+y) + v_2(n) - 1$ for even $n$ (so $v_2(x-y) + v_2(n)$ if $4 \mid x - y$).

% \section{Fractions}
	% \kactlimport{ContinuedFractions.h}
	% \kactlimport{FracBinarySearch.h}

\section{Pythagorean Triples}
 The Pythagorean triples are uniquely generated by
 \[ a=k\cdot (m^{2}-n^{2}),\ \,b=k\cdot (2mn),\ \,c=k\cdot (m^{2}+n^{2}), \]
 with $m > n > 0$, $k > 0$, $m \bot n$, and either $m$ or $n$ even.

\section{Primes}
	$p=962592769$ is such that $2^{21} \mid p-1$, which may be useful. For hashing
	use 970592641 (31-bit number), 31443539979727 (45-bit), 3006703054056749
	(52-bit). There are 78498 primes less than 1\,000\,000.

	Primitive roots exist modulo any prime power $p^a$, except for $p = 2, a > 2$, and there are $\phi(\phi(p^a))$ many.
	For $p = 2, a > 2$, the group $\mathbb Z_{2^a}^\times$ is instead isomorphic to $\mathbb Z_2 \times \mathbb Z_{2^{a-2}}$.

\section{Floor blocks}
	$\lfloor n/i \rfloor$ takes $O(\sqrt n)$ distinct values; the blocks $[l, r]$ with equal $\lfloor n/i \rfloor$ are:\\
	\verb@for (l = 1; l <= n; l = r + 1) {@\\
	\verb@  r = n / (n / l); ... }@\\
	$\lfloor \lfloor n/a \rfloor / b \rfloor = \lfloor n/(ab) \rfloor$, and
	$\sum_{i=1}^{n} \lfloor n/i \rfloor = 2 \sum_{i \le \lfloor\sqrt n\rfloor} \lfloor n/i \rfloor - \lfloor\sqrt n\rfloor^2$.
	\kactlimport{Lucy.h}
	\kactlimport{DuSieve.h}

\section{Estimates}
	$\sum_{d|n} d = O(n \log \log n)$.

	The number of divisors of $n$ is at most around 100 for $n < 5e4$, 500 for $n < 1e7$, 2000 for $n < 1e10$, 200\,000 for $n < 1e19$.

\section{Mobius Function}
\[
	\mu(n) = \begin{cases} 0 & n \textrm{ is not square free}\\ 1 & n \textrm{ has even number of prime factors}\\ -1 & n \textrm{ has odd number of prime factors}\\\end{cases}
\]
  Mobius Inversion:
  \[ g(n) = \sum_{d|n} f(d) \Leftrightarrow f(n) = \sum_{d|n} \mu(d)g(n/d) \]
  Other useful formulas/forms:

  $ \sum_{d | n} \mu(d) = [ n = 1] $ (very useful)

  $ g(n) = \sum_{n|d} f(d) \Leftrightarrow f(n) = \sum_{n|d} \mu(d/n)g(d)$

 $ g(n) = \sum_{1 \leq m \leq n} f(\left\lfloor\frac{n}{m}\right \rfloor ) \Leftrightarrow f(n) = \sum_{1\leq m\leq n} \mu(m)g(\left\lfloor\frac{n}{m}\right\rfloor)$
